\documentclass[11pt,a4paper]{article}
\usepackage[a4paper,margin=1in]{geometry}
\usepackage[T1]{fontenc}
\usepackage[utf8]{inputenc}
\usepackage{amsmath,amssymb,amsthm,bm,mathtools}
\usepackage{booktabs}
\usepackage{enumitem}
\usepackage{xcolor}
\usepackage{hyperref}
\hypersetup{colorlinks=true,linkcolor=blue!50!black,citecolor=blue!50!black,urlcolor=blue!50!black}

\newtheorem{theorem}{Theorem}
\newtheorem{definition}{Definition}
\newtheorem{remark}{Remark}

\newcommand{\bbC}{\mathbb{C}}
\newcommand{\bbI}{\mathbb{I}}
\newcommand{\cH}{\mathcal{H}}
\newcommand{\Graph}{G}
\newcommand{\Gammaop}{\Gamma}
\newcommand{\Pgraph}{\mathfrak{P}_{\mathrm{graph}}}
\newcommand{\kerdim}{\dim\ker}
\providecommand{\doi}[1]{\href{https://doi.org/#1}{doi:#1}}

\title{Chiral balance without Brillouin:\\
graph-local constraints on emergent zero-mode sectors}
\author{Nelson Bolivar\\
\small Astrum Drive Technologies\\
\small \href{mailto:astrumdrivetechnologies@gmail.com}{astrumdrivetechnologies@gmail.com}}
\date{Draft 1, May 2026}

\begin{document}
\maketitle

\begin{abstract}
We study a structural version of chiral balance that does not rely on Bloch
periodicity or a Brillouin torus. The usual lattice discussion of Weyl and
Dirac fermions is naturally formulated in reciprocal space, where nodes,
chiral charges, and doubling constraints are expressed over a compact momentum
manifold. This formulation is powerful for periodic lattices but is not
obviously primitive for a more general discrete substrate. We replace the
periodic lattice by finite graph-local data: a graph of sites or events, a
finite-dimensional internal fiber, a local Hamiltonian or discrete-time step,
and, when present, a chiral grading. We review the elementary spectral pairing
implied by $\{\Gamma,H\}=0$, then sharpen the finite-graph statement in two
directions. First, for generic weighted bipartite graph Hamiltonians, the
zero-mode count is controlled by the maximum matching number rather than only
by sublattice imbalance:
\[
  \dim\ker H = |A|+|B|-2\nu(G).
\]
Second, if the grading is only approximate,
$\|\{\Gamma,H\}\|\le \epsilon$, exact pairing is replaced by an
$\epsilon$-pairing of the spectrum, and the zero band of the exactly chiral
part remains inside an $\epsilon/2$ window. These results do not constitute a
Nielsen--Ninomiya theorem beyond periodic lattices. They isolate a real-space
operator layer beneath it: chiral balance can be formulated as matching,
locality, and approximate-grading constraints on finite graph-local operators,
without using reciprocal-space periodicity.
\end{abstract}

\section{Introduction}

The first paper in this sequence showed how Dirac covariance can emerge as an
infrared organization of local discrete dynamics. The present paper asks a
deeper structural question: which parts of the Dirac/Weyl obstruction survive
when reciprocal space is no longer assumed to be fundamental?

The usual route to lattice fermions begins with a periodic lattice. Fourier
transform converts locality into a smooth matrix-valued function $H(k)$ over a
Brillouin torus. Weyl nodes become zeros of this function, chirality is read
from local orientation data, and doubling constraints become global
cancellation statements over the compact reciprocal manifold. This is the
natural language of the Nielsen--Ninomiya theorem
\cite{NielsenNinomiya1981a,NielsenNinomiya1981b}.

The motivating question here is narrower and deliberately safer:
\begin{quote}
Can locality, chirality, and Dirac-sector constraints be formulated directly
on a discrete substrate, without passing first through Bloch momentum?
\end{quote}
We do not claim to prove a full Nielsen--Ninomiya theorem beyond periodic
lattices. Instead, we isolate a finite-graph analogue of chiral balance:
spectral pairing, zero-mode constraints, and index-type obstructions can be
formulated without reciprocal-space periodicity, provided the microscopic
dynamics carries a local graph structure and an internal chiral grading.

\section{From Bloch Space to Graph-Local Data}

The Bloch representation is useful when translation symmetry exists, but it is
not ontologically neutral. It assumes that the microscopic substrate admits a
global reciprocal description. If the substrate is a finite graph, a
causal-local network, or a locally regular but globally non-periodic structure,
then reciprocal space may not exist as primitive data.

We therefore separate two layers.
\begin{enumerate}[label=(\alph*)]
\item The \emph{reciprocal/topological layer}, where Weyl points are zeros of
      $H(k)$ and total chirality is a global degree or cancellation condition.
\item The \emph{structural/local layer}, where sites, internal fibers,
      locality, and algebraic gradings are defined directly.
\end{enumerate}
This draft studies the second layer.

\section{Minimal Graph-Local Protospace}

\begin{definition}[Graph-local protospace]
A minimal graph-local protospace candidate is a tuple
\[
  \Pgraph = (\Graph,F,\cH,\mathcal{D},\Gammaop,\mathcal{G}_{\mathrm{micro}}),
\]
where $\Graph=(V,E)$ is a finite graph or a family of finite graphs, $F$ is a
finite-dimensional internal fiber, $\cH=\ell^2(V)\otimes F$ is the microscopic
Hilbert space, $\mathcal{D}$ is either a Hamiltonian $H$ or a discrete-time
unitary update $U$, $\Gammaop$ is a chiral grading when present, and
$\mathcal{G}_{\mathrm{micro}}$ denotes exact microscopic symmetries.
\end{definition}

The key move is that reciprocal space is not included as primitive data. It may
emerge, or become a useful representation, when $\Graph$ is periodic. The
primitive object is a local operator over a graph with finite internal
structure.

\section{Locality Without Fourier Transform}

\begin{definition}[$R$-local Hamiltonian]
Let $d_G$ be the graph distance on $G$. A Hamiltonian $H$ on
$\ell^2(V)\otimes F$ is $R$-local if
\[
  \langle x,a|H|y,b\rangle = 0
  \qquad \text{whenever } d_G(x,y)>R,
\]
for vertices $x,y\in V$ and internal labels $a,b\in F$.
\end{definition}

For a discrete-time update $U$, the same definition applies to the matrix
elements of $U$. This is enough to formulate microscopic finite propagation
speed and local causal neighborhoods without invoking reciprocal space.

\section{Chiral Grading and Spectral Pairing}

Let $H$ be a finite-dimensional Hermitian operator and let $\Gammaop$ be a
unitary involution:
\[
  \Gammaop^2=\bbI,\qquad \Gammaop^\dagger=\Gammaop.
\]
We say that $H$ is chiral with respect to $\Gammaop$ if
\[
  \{\Gammaop,H\}=\Gammaop H+H\Gammaop=0.
\]
This is the finite graph-local form of the chiral grading familiar from the
Altland--Zirnbauer symmetry classes \cite{AltlandZirnbauer1997}.

\begin{theorem}[Spectral pairing]
If $H\psi=E\psi$ and $E\neq 0$, then
\[
  H(\Gammaop\psi)=-E(\Gammaop\psi).
\]
Consequently every nonzero eigenvalue appears in an opposite-sign pair.
\end{theorem}

\begin{proof}
Using $\Gammaop H=-H\Gammaop$,
\[
  H\Gammaop\psi=-\Gammaop H\psi=-E\Gammaop\psi.
\]
Thus $\Gammaop\psi$ is an eigenvector of $H$ with eigenvalue $-E$.
\end{proof}

The proof is purely algebraic. It does not use translation symmetry, Fourier
transform, Bloch momentum, or a Brillouin torus.

\section{Bipartite Graphs and Zero-Mode Constraints}

Let $G$ be bipartite with $V=A\sqcup B$ and edges only between $A$ and $B$. A
nearest-neighbor hopping Hamiltonian has block form
\[
  H =
  \begin{pmatrix}
    0 & T\\
    T^\dagger & 0
  \end{pmatrix},
\]
where $T:\bbC^{|B|}\to\bbC^{|A|}$. The natural sublattice grading is
\[
  \Gammaop =
  \begin{pmatrix}
    \bbI_A & 0\\
    0 & -\bbI_B
  \end{pmatrix}.
\]
Then $\Gammaop H\Gammaop=-H$.

\begin{theorem}[Bipartite zero-mode bound]
For the bipartite Hamiltonian above,
\[
  \kerdim H \geq \big||A|-|B|\big|.
\]
\end{theorem}

\begin{proof}
The kernel of $H$ consists of pairs $(\psi_A,\psi_B)$ satisfying
\[
  T\psi_B=0,\qquad T^\dagger\psi_A=0.
\]
Thus
\[
  \kerdim H = \dim\ker T+\dim\ker T^\dagger.
\]
By rank-nullity,
\[
  \dim\ker T = |B|-\operatorname{rank}T,\qquad
  \dim\ker T^\dagger = |A|-\operatorname{rank}T.
\]
Therefore
\[
  \kerdim H=|A|+|B|-2\operatorname{rank}T.
\]
Since $\operatorname{rank}T\leq \min(|A|,|B|)$, the claimed lower bound follows.
\end{proof}

Sublattice imbalance forces zero modes. This is an index-type statement on
finite graphs and does not require momentum space. In condensed-matter form,
closely related imbalance statements appear in Lieb's theorem for bipartite
Hubbard systems \cite{Lieb1989}; in graph-theoretic language, the sharper
object is not only the imbalance but the maximum matching structure of the
underlying bipartite graph \cite{DulmageMendelsohn1958}.
For disordered bipartite hopping models, the same off-diagonal structure is
also the source of robust zero-energy phenomena \cite{InuiTrugmanAbrahams1994}.

\section{Matching-Refined Chiral Nullity}

The imbalance bound is universal but not sharp. A graph can have
$|A|=|B|$ and still support zero modes because its connectivity fails to admit
a perfect matching. This gives a more precise real-space replacement for the
first layer of chiral counting.

Let $\nu(G)$ denote the size of a maximum matching of the bipartite graph
$G=(A\sqcup B,E)$. Let $T(w)$ be the $|A|\times |B|$ weighted bi-adjacency
matrix with one independent nonzero weight on each edge and zero entries on
non-edges.

\begin{theorem}[Generic matching nullity]
For generic edge weights on a finite bipartite graph,
\[
  \operatorname{rank} T(w) = \nu(G).
\]
Consequently the off-diagonal chiral Hamiltonian
\[
  H(w)=
  \begin{pmatrix}
    0 & T(w)\\
    T(w)^\dagger & 0
  \end{pmatrix}
\]
has generic nullity
\[
  \dim\ker H(w)=|A|+|B|-2\nu(G).
\]
\end{theorem}

\begin{proof}
Every nonzero $r\times r$ minor of $T(w)$ expands as a sum over matchings of
size $r$ in the corresponding induced bipartite subgraph. Therefore no minor
of size larger than $\nu(G)$ can be nonzero. Conversely, a matching of size
$\nu(G)$ selects a square submatrix whose determinant contains a monomial equal
to the product of the matched edge weights. With algebraically independent
edge weights this determinant polynomial is not identically zero, so the
generic rank is at least $\nu(G)$. The two inequalities give
$\operatorname{rank}T(w)=\nu(G)$ for all weights outside the algebraic
exceptional set where the relevant minors accidentally vanish. Since
\[
  \dim\ker H = \dim\ker T+\dim\ker T^\dagger
             = |B|-\operatorname{rank}T + |A|-\operatorname{rank}T,
\]
the nullity formula follows.
\end{proof}

This theorem is not a claim that the graph theorem itself is new. Rather, its
role here is to identify the real-space combinatorial datum that replaces the
first pass of Brillouin-zone node counting when translation symmetry is not
primitive. The sublattice-imbalance bound is the coarse corollary
\[
  |A|+|B|-2\nu(G)\geq \big||A|-|B|\big|.
\]

\section{Approximate Chiral Balance}

Exact microscopic chiral symmetry is often too rigid for an emergent theory.
The useful question is what survives if the grading is only approximate.

\begin{theorem}[Approximate spectral pairing]
Let $H$ be finite-dimensional Hermitian and let $\Gamma$ be a unitary
involution. If
\[
  \|\{\Gamma,H\}\|\leq \epsilon,
\]
then for every eigenvalue $E\in\sigma(H)$ there exists
$E'\in\sigma(H)$ such that
\[
  |E+E'|\leq \epsilon.
\]
\end{theorem}

\begin{proof}
Let $H\psi=E\psi$ with $\|\psi\|=1$, and set $\phi=\Gamma\psi$. Then
\[
  (H+E\bbI)\phi
  = (H\Gamma+E\Gamma)\psi
  = (\{\Gamma,H\}-\Gamma H)\psi+E\Gamma\psi
  = \{\Gamma,H\}\psi.
\]
Hence $\|(H-(-E)\bbI)\phi\|\leq \epsilon$. For a Hermitian matrix, a
normalized vector with residual at most $\epsilon$ from a real number
$\lambda$ implies that the spectrum has an eigenvalue within $\epsilon$ of
$\lambda$. Taking $\lambda=-E$ gives the result.
\end{proof}

There is also a useful zero-band version. Decompose
\[
  H_{-}=\frac{H-\Gamma H\Gamma}{2},\qquad
  H_{+}=\frac{H+\Gamma H\Gamma}{2}.
\]
Then $\{\Gamma,H_-\}=0$, while $H_+$ is the grading-breaking part and
$\|H_+\|\leq \epsilon/2$. If $H_-$ has $d$ exact structural zero modes, Weyl
stability implies that $H=H_-+H_+$ has at least $d$ eigenvalues in
$[-\epsilon/2,\epsilon/2]$. Thus approximate chirality converts exact
structural zeros into a controlled small mass window rather than destroying
the low-energy sector outright.

\section{A Local Boundary Obstruction}

As a minimal illustration of what a real-space obstruction looks like, consider
the chain segment
\[
  A_1 - B_1 - A_2 - B_2 - A_3
\]
with hopping weights $t_1,t_2,t_3,t_4$. A zero mode supported only on the
interior site $A_2$ would have $\psi_{A_1}=\psi_{A_3}=0$ and
$\psi_{A_2}\neq0$. The zero-energy equations at $B_1$ and $B_2$ reduce to
\[
  t_2\psi_{A_2}=0,\qquad t_3\psi_{A_2}=0.
\]
Thus a one-site isolated bulk zero mode forces $t_2=t_3=0$, i.e. it cuts the
site away from the connected graph.

This is only a local lemma, not a general no-go theorem for all compact
localized states: on more complicated graphs, multi-site destructive
interference can produce compact zero modes without severing every adjacent
edge. Its value here is diagnostic. A credible graph-local replacement for
reciprocal-space doubling must specify which kinds of localization are
structural, which are fine-tuned, and which require actual graph cuts.

\section{Relation to Nielsen--Ninomiya}

The Nielsen--Ninomiya theorem is a reciprocal-space/topological obstruction for
chiral fermions on lattices under assumptions such as locality, Hermiticity,
and translation invariance. In the usual periodic setting, Weyl points carry
chiral charges and the total charge over the Brillouin torus cancels.

The present paper does not replace that theorem. It separates the obstruction
into two layers:
\begin{description}[leftmargin=2.6cm,style=nextline]
\item[Layer A] Weyl nodes are zeros of $H(k)$; chirality is a local degree of
a map; global cancellation follows from compactness of the Brillouin zone.
\item[Layer B] Finite graph-local operators with chiral grading exhibit
spectral pairing, matching-refined zero-mode constraints, and approximate
pairing bounds, independently of reciprocal space.
\end{description}
The claim is that Layer B survives beyond Bloch periodicity. It is weaker than
the full reciprocal/topological theorem, but stronger than the elementary
observation that a block off-diagonal matrix has a symmetric spectrum.

\section{Why This Matters for Protospace}

If protospace is defined only as a periodic lattice, then Bloch space and
Brillouin-zone topology dominate the analysis. If protospace is a more general
local discrete substrate, the first question is not ``what is $H(k)$?'' but:
\begin{quote}
What local operator structure is sufficient to support an infrared Dirac
sector?
\end{quote}

The results above suggest that the minimum ingredients include local
connectivity, a finite internal fiber, a grading or approximate grading, a
mechanism for infrared linearization, stability of the low-energy sector under
local perturbations, and a structural version of chirality balance. A genuine
Dirac sector still requires more than the results proved here: one needs an
infrared notion of momentum, a linear dispersion law, or an effective continuum
limit.

\section{Verified Graph Families}

The repository contains a first executable validation layer for this draft in
\path{validators/graph_chiral_balance.py} and
\path{tests/test_graph_chiral_balance.py}. The checks are deliberately
finite-dimensional and do not invoke Fourier transform.

\begin{center}
\small
\begin{tabular}{p{0.24\linewidth}p{0.27\linewidth}p{0.38\linewidth}}
\toprule
Claim & Mathematical object & Validator \\
\midrule
Matching nullity & generic bipartite graph & \path{test_generic_matching_nullity} \\
Approximate pairing & $\epsilon$-chiral finite $H$ & \path{test_approximate_pairing_bound} \\
Zero-band stability & perturbed imbalanced graph & \path{test_structural_zero_modes_stay_in_epsilon_window} \\
Local obstruction & symbolic chain segment & \path{test_single_site_bulk_mode_requires_edge_cut} \\
Spectral pairing & finite chiral $H$ & \path{test_chiral_pairing_finite_matrix} \\
Bipartite index bound & bipartite blocks & \path{test_bipartite_index_bound} \\
SSH edge zero modes & open SSH chain & \path{test_ssh_open_chain_zero_modes} \\
Chiral disorder & disordered bipartite hopping & \path{test_chiral_disorder_preserves_pairing} \\
Chiral breaking & onsite perturbation & \path{test_chiral_breaking_lifts_pairing} \\
Graph locality & irregular finite graph & \path{test_irregular_graph_locality} \\
\bottomrule
\end{tabular}
\end{center}

\section{Relation to Discrete-Time Quantum Walks}

For Hamiltonian systems, chiral symmetry is naturally expressed around
$E=0$. For discrete-time unitary walks, the spectrum is quasienergy:
\[
  U\psi = e^{-i\omega\Delta t}\psi.
\]
Quasienergy is periodic, so special sectors occur near $\omega=0$ and
$\omega=\pi/\Delta t$ \cite{Asboth2012}. The graph-local Hamiltonian results developed here
should therefore be generalized to unitary chiral walks, where the relevant
spectral obstructions live on the quasienergy circle rather than on the real
line. We leave this as a controlled next extension.

\section{Open Problems}

\begin{description}[leftmargin=1.2cm,style=nextline]
\item[E1] \textbf{Dirac linearization without translation symmetry.}
Spectral pairing and zero-mode constraints are graph-local. A genuine Dirac
sector also requires an infrared notion of momentum, linear dispersion, or an
effective continuum limit.

\item[E2] \textbf{Chiral balance for unitary quantum walks.}
In discrete time, quasienergy is circular. The analogue of chiral balance must
account for both $0$ and $\pi$ quasienergy sectors.

\item[E3] \textbf{Structural Nielsen--Ninomiya.}
Can one formulate a no-go theorem directly in terms of graph-locality, grading,
and finite propagation, without assuming a Brillouin torus? This draft provides
algebraic ingredients but not the full theorem.
\end{description}

\section{Conclusion}

This paper isolates a structural layer beneath the usual reciprocal-space
discussion of chiral lattice fermions. Chiral spectral pairing, bipartite
zero-mode constraints, matching-refined generic nullity, and
$\epsilon$-approximate pairing can be formulated directly on finite
graph-local operators. These statements require neither Bloch waves nor a
Brillouin torus.

The result is not a replacement for Nielsen--Ninomiya. Rather, it identifies
which part of the obstruction survives when reciprocal space is removed from
the primitive data. The matching number $\nu(G)$ becomes the finite real-space
datum controlling generic chiral nullity, while approximate grading bounds
show that small symmetry leakage produces controlled small mass windows rather
than arbitrary spectral collapse. This gives a sharper meaning to protospace:
a candidate microscopic substrate is not merely a discrete set of sites, but a
graph-local operator system with internal fiber, locality, grading structure,
matching constraints, and controlled infrared sectors.

\section*{Code, Data, and AI Assistance}

The validation code for this draft is maintained in the companion repository
\url{https://github.com/xys004/protoespacio-validation}. A Zenodo DOI should be
minted for a frozen Paper II release once the present draft and its validators
are stable.

AI-assisted coding and editorial tools were used to help draft, refactor, and
review parts of the validation modules, tests, and manuscript text. All
scientific claims, validator design choices, and final interpretations remain
the responsibility of the author.

\begin{thebibliography}{9}

\bibitem{NielsenNinomiya1981a}
H.~B. Nielsen and M.~Ninomiya.
\newblock Absence of neutrinos on a lattice. I. Proof by homotopy theory.
\newblock \emph{Nuclear Physics B}, 185(1):20--40, 1981.

\bibitem{NielsenNinomiya1981b}
H.~B. Nielsen and M.~Ninomiya.
\newblock Absence of neutrinos on a lattice. II. Intuitive topological proof.
\newblock \emph{Nuclear Physics B}, 193(1):173--194, 1981.

\bibitem{AltlandZirnbauer1997}
A.~Altland and M.~R. Zirnbauer.
\newblock Nonstandard symmetry classes in mesoscopic normal-superconducting
hybrid structures.
\newblock \emph{Physical Review B}, 55:1142--1161, 1997.
\newblock \doi{10.1103/PhysRevB.55.1142}.

\bibitem{Lieb1989}
E.~H. Lieb.
\newblock Two theorems on the Hubbard model.
\newblock \emph{Physical Review Letters}, 62(10):1201--1204, 1989.
\newblock \doi{10.1103/PhysRevLett.62.1201}.

\bibitem{DulmageMendelsohn1958}
A.~L. Dulmage and N.~S. Mendelsohn.
\newblock Coverings of bipartite graphs.
\newblock \emph{Canadian Journal of Mathematics}, 10:517--534, 1958.
\newblock \doi{10.4153/CJM-1958-052-0}.

\bibitem{InuiTrugmanAbrahams1994}
M.~Inui, S.~A. Trugman, and E.~Abrahams.
\newblock Unusual properties of midband states in systems with off-diagonal
disorder.
\newblock \emph{Physical Review B}, 49:3190--3196, 1994.
\newblock \doi{10.1103/PhysRevB.49.3190}.

\bibitem{Asboth2012}
J.~K. Asb{\'o}th.
\newblock Symmetries, topological phases, and bound states in the
one-dimensional quantum walk.
\newblock \emph{Physical Review B}, 86:195414, 2012.
\newblock \doi{10.1103/PhysRevB.86.195414}.

\end{thebibliography}

\end{document}
